% GENERATED FILE. Edit canonical lessons, then run npm run generate:books.
% canonical-source-hash: fnv1a64:48545a83c07d837a
% canonical-lessons: HI-C515-dhila, HI-C515-penc, HI-C515-peckas, HI-C515-kasna, HI-C515-kharkharana

\chapter{Something Has Come Loose}
\label{ch:hi-something-has-come-loose}
\hlchaptermodality{515}

\begin{chapteropening}
\textbf{By the end of this chapter:} \emph{I can say what is loose and fix it.}

Complete the last lesson of chapter 515: I can say what is loose and fix it.
\end{chapteropening}

\section[\texorpdfstring{\hi{ढीला} (ḍhīlā)}{ढीला (ḍhīlā)}]{\hi{ढीला} (ḍhīlā) — loose; slack}

\label{lesson:HI-C515-dhila}



\begin{quote}
Before the new one: say the Hindi for to polish, to make shine, then the Hindi for to melt (something).
\end{quote}

\subsection*{You'll want to know: \hi{ढीला}}
\textbf{\hi{ढीला}} — \emph{ḍhīlā} — \textquotedblleft{}loose; slack\textquotedblright{}.

\begin{quote}
\hi{यह} \hi{कमीज़} \hi{मुझे} \hi{बहुत} \hi{ढीली} \hi{है।}
\end{quote}

\begin{quote}
\emph{Yah kamīz mujhe bahut ḍhīlī hai.}
\end{quote}

\begin{quote}
This shirt is very loose on me.
\end{quote}

\begin{cousinweb}
\textbf{\hi{ढीला}} (\emph{ḍhīlā}) is a tadbhava, inherited from \emph{*ḍhila}, \textquotedblleft{}slack\textquotedblright{}, a form CDIAL reconstructs from its descendants (CDIAL 5590).
\end{cousinweb}

\begin{grammarlens}[title={What you've built}]
One more word for describing what you lived and felt.
\end{grammarlens}

\subsection*{Your turn}
\begin{itemize}
\raggedright
  \item \emph{Say it:} \emph{ḍhīlā}
  \item \emph{Say it:} \emph{ḍhīlā}, once more
  \item \emph{Say it:} \emph{Yah kamīz mujhe bahut ḍhīlī hai.} — This shirt is very loose on me.
  \item \emph{Recall:} say the Hindi for to polish, to make shine, then the Hindi for to melt (something), then say \emph{ḍhīlā} again
\end{itemize}

\begin{tcolorbox}[breakable,colback=teal!4,colframe=teal!35!black,title={Before you move on}]
What does \hi{ढीला} mean? (\textquotedblleft{}Loose; slack\textquotedblright{}.) Say it once more.
\end{tcolorbox}
\section[\texorpdfstring{\hi{पेंच} (peñc)}{पेंच (peñc)}]{\hi{पेंच} (peñc) — a screw}

\label{lesson:HI-C515-penc}



\begin{quote}
Before the new one: say the Hindi for to melt (something), then the Hindi for loose.
\end{quote}

\subsection*{You'll want to know: \hi{पेंच}}
\textbf{\hi{पेंच}} — \emph{peñc} — \textquotedblleft{}a screw\textquotedblright{}.

\begin{quote}
\hi{पंखे} \hi{का} \hi{एक} \hi{पेंच} \hi{ढीला} \hi{हो} \hi{गया।}
\end{quote}

\begin{quote}
\emph{Paṅkhe kā ek peñc ḍhīlā ho gayā.}
\end{quote}

\begin{quote}
One of the fan's screws has come loose.
\end{quote}

Masculine.

\begin{cousinweb}
\textbf{\hi{पेंच}} (\emph{peñc}) is usually traced to Persian \emph{pēch}, \textquotedblleft{}a twist, a screw\textquotedblright{}.
\end{cousinweb}

\begin{grammarlens}[title={What you've built}]
One more word for describing what you lived and felt.
\end{grammarlens}

\subsection*{Your turn}
\begin{itemize}
\raggedright
  \item \emph{Say it:} \emph{peñc}
  \item \emph{Say it:} \emph{peñc}, once more
  \item \emph{Say it:} \emph{Paṅkhe kā ek peñc ḍhīlā ho gayā.} — One of the fan's screws has come loose.
  \item \emph{Recall:} say the Hindi for to melt (something), then the Hindi for loose, then say \emph{peñc} again
\end{itemize}

\begin{tcolorbox}[breakable,colback=teal!4,colframe=teal!35!black,title={Before you move on}]
What does \hi{पेंच} mean? (\textquotedblleft{}A screw\textquotedblright{}.) Say it once more.
\end{tcolorbox}
\section[\texorpdfstring{\hi{पेचकस} (peckas)}{पेचकस (peckas)}]{\hi{पेचकस} (peckas) — a screwdriver}

\label{lesson:HI-C515-peckas}



\begin{quote}
Before the new one: say the Hindi for loose, then the Hindi for a screw.
\end{quote}

\subsection*{You'll want to know: \hi{पेचकस}}
\textbf{\hi{पेचकस}} — \emph{peckas} — \textquotedblleft{}a screwdriver\textquotedblright{}.

\begin{quote}
\hi{पेचकस} \hi{अलमारी} \hi{में} \hi{ऊपर} \hi{रखा} \hi{है।}
\end{quote}

\begin{quote}
\emph{Peckas almārī meṁ ūpar rakhā hai.}
\end{quote}

\begin{quote}
The screwdriver is up in the cupboard.
\end{quote}

Masculine.

\begin{cousinweb}
\textbf{\hi{पेचकस}} (\emph{peckas}) is usually traced to Persian \emph{pēch}, \textquotedblleft{}a screw\textquotedblright{}, and \emph{kash}, \textquotedblleft{}pulling, drawing\textquotedblright{}: the thing that turns screws.
\end{cousinweb}

\begin{grammarlens}[title={What you've built}]
One more word for describing what you lived and felt.
\end{grammarlens}

\subsection*{Your turn}
\begin{itemize}
\raggedright
  \item \emph{Say it:} \emph{peckas}
  \item \emph{Say it:} \emph{peckas}, once more
  \item \emph{Say it:} \emph{Peckas almārī meṁ ūpar rakhā hai.} — The screwdriver is up in the cupboard.
  \item \emph{Recall:} say the Hindi for loose, then the Hindi for a screw, then say \emph{peckas} again
\end{itemize}

\begin{tcolorbox}[breakable,colback=teal!4,colframe=teal!35!black,title={Before you move on}]
What does \hi{पेचकस} mean? (\textquotedblleft{}A screwdriver\textquotedblright{}.) Say it once more.
\end{tcolorbox}
\section[\texorpdfstring{\hi{कसना} (kasnā)}{कसना (kasnā)}]{\hi{कसना} (kasnā) — to tighten}

\label{lesson:HI-C515-kasna}



\begin{quote}
Before the new one: say the Hindi for a screw, then the Hindi for a screwdriver.
\end{quote}

\subsection*{You'll want to know: \hi{कसना}}
\textbf{\hi{कसना}} — \emph{kasnā} — \textquotedblleft{}to tighten\textquotedblright{}.

\begin{quote}
\hi{अरुण} \hi{ने} \hi{पेचकस} \hi{से} \hi{सब} \hi{पेंच} \hi{कस} \hi{दिए।}
\end{quote}

\begin{quote}
\emph{Aruṇ ne peckas se sab peñc kas diye.}
\end{quote}

\begin{quote}
Arun tightened all the screws with the screwdriver.
\end{quote}

Past, with \hi{ने} (he or she …): \hi{उसने} \hi{पेंच} \hi{कसा}, \emph{usne peñc kasā}.

\begin{cousinweb}
\textbf{\hi{कसना}} (\emph{kasnā}) is a tadbhava, inherited from Sanskrit \emph{karṣati}, \textquotedblleft{}draws, pulls\textquotedblright{} (CDIAL 2908). CDIAL gives the Hindi verb as \textquotedblleft{}to tighten\textquotedblright{}.
\end{cousinweb}

\begin{grammarlens}[title={What you've built}]
One more word for describing what you lived and felt.
\end{grammarlens}

\subsection*{Your turn}
\begin{itemize}
\raggedright
  \item \emph{Say it:} \emph{kasnā}
  \item \emph{Say it:} \emph{kasnā}, once more
  \item \emph{Say it:} \emph{Aruṇ ne peckas se sab peñc kas diye.} — Arun tightened all the screws with the screwdriver.
  \item \emph{Recall:} say the Hindi for a screw, then the Hindi for a screwdriver, then say \emph{kasnā} again
\end{itemize}

\begin{tcolorbox}[breakable,colback=teal!4,colframe=teal!35!black,title={Before you move on}]
What does \hi{कसना} mean? (\textquotedblleft{}To tighten\textquotedblright{}.) Say it once more.
\end{tcolorbox}
\section[\texorpdfstring{\hi{खड़खड़ाना} (khaṛkhaṛānā)}{खड़खड़ाना (khaṛkhaṛānā)}]{\hi{खड़खड़ाना} (khaṛkhaṛānā) — to rattle}

\label{lesson:HI-C515-kharkharana}



\begin{quote}
Before the new one: say the Hindi for a screwdriver, then the Hindi for to tighten.
\end{quote}

\subsection*{You'll want to know: \hi{खड़खड़ाना}}
\textbf{\hi{खड़खड़ाना}} — \emph{khaṛkhaṛānā} — \textquotedblleft{}to rattle\textquotedblright{}.

\begin{quote}
\hi{हवा} \hi{में} \hi{पुरानी} \hi{खिड़की} \hi{खड़खड़ा} \hi{रही} \hi{है।}
\end{quote}

\begin{quote}
\emph{Havā meṁ purānī khiṛkī khaṛkhaṛā rahī hai.}
\end{quote}

\begin{quote}
The old window is rattling in the wind.
\end{quote}

Past (it …): \hi{वह} \hi{दरवाज़ा} \hi{खड़खड़ाया}, \emph{vah darvāzā khaṛkhaṛāyā}.

\begin{cousinweb}
\textbf{\hi{खड़खड़ाना}} (\emph{khaṛkhaṛānā}) is a tadbhava, inherited from Sanskrit \emph{khaṭakhaṭāyatē}, \textquotedblleft{}crackles\textquotedblright{}, a sound word (CDIAL 3771). CDIAL gives the Hindi verb as \textquotedblleft{}to make a noise with\textquotedblright{}.
\end{cousinweb}

\begin{grammarlens}[title={What you've built}]
One more word for describing what you lived and felt.
\end{grammarlens}

\subsection*{Your turn}
\begin{itemize}
\raggedright
  \item \emph{Say it:} \emph{khaṛkhaṛānā}
  \item \emph{Say it:} \emph{khaṛkhaṛānā}, once more
  \item \emph{Say it:} \emph{Havā meṁ purānī khiṛkī khaṛkhaṛā rahī hai.} — The old window is rattling in the wind.
  \item \emph{Recall:} say the Hindi for a screwdriver, then the Hindi for to tighten, then say \emph{khaṛkhaṛānā} again
\end{itemize}

\begin{tcolorbox}[breakable,colback=teal!4,colframe=teal!35!black,title={Before you move on}]
What does \hi{खड़खड़ाना} mean? (\textquotedblleft{}To rattle\textquotedblright{}.) Say it once more.
\end{tcolorbox}
